\section{Empirical Application}\label{sec:application}

% ---------------------------------------------------------------------------
\subsection{Emergency-department visits across Italian municipalities}\label{sec:app_data}

To illustrate the ZISAR model in a substantively relevant setting, we analyse the spatial distribution of emergency-department (ED) visit counts across municipalities in the Piedmont region of north-western Italy.
ED utilisation is a canonical application for zero-inflated spatial count models: many small municipalities record zero visits for specific conditions during a given period (structural zeros arising from the absence of at-risk populations, or from residents seeking care in a neighbouring municipality's hospital); the non-zero counts exhibit substantial overdispersion driven by heterogeneity in population size, age structure, and access to primary care; and neighbouring municipalities share health-system characteristics and patient-flow patterns that induce spatial dependence \citep{neelon2013,baldi2014}.

\paragraph{Data.}
We use municipality-level data for $n = 300$ municipalities in the Piedmont region.
The outcome variable $y_i$ is the total number of ED visits (all causes) recorded in municipality $i$ during the observation period.
The data exhibit clear signatures of both zero-inflation and overdispersion: approximately 25\% of municipalities report zero visits, while the sample variance substantially exceeds the sample mean.

\paragraph{Covariates.}
Five municipality-level covariates enter the count process ($\bX$):
\begin{itemize}
  \item \texttt{pop}: log-population (standardised);
  \item \texttt{aging}: aging index---ratio of the population aged 65+ to the population aged 0--14 (standardised);
  \item \texttt{dist\_hosp}: distance to the nearest hospital with an ED (km, standardised);
  \item \texttt{income}: median declared income per capita (standardised).
\end{itemize}
Two covariates enter the zero-inflation process ($\bZ$):
\begin{itemize}
  \item \texttt{dist\_hosp}: distance to the nearest hospital (standardised);
  \item \texttt{gp\_density}: density of general practitioners per 1{,}000 residents (standardised).
\end{itemize}
The rationale is that the probability of a \emph{structural} zero---reflecting the absence of any demand, as opposed to a sampling zero from low intensity---depends primarily on geographic accessibility (distance to hospital) and the availability of primary-care alternatives (GP density) that may absorb demand.

\paragraph{Spatial weights.}
The spatial weights matrix $\bW$ is constructed from 6-nearest-neighbour adjacency (based on Euclidean distance between municipality centroids) and row-standardised.
Sensitivity to alternative choices ($k \in \{4, 8\}$ and queen contiguity) is examined in Appendix~\ref{app:sensitivity}.

% ---------------------------------------------------------------------------
\subsection{Results}\label{sec:app_results}

We estimate three specifications:
\begin{enumerate}
  \item \textbf{ZINB}: Standard zero-inflated NB, ignoring spatial dependence.
  \item \textbf{SAR-NB}: NB GLM with an instrumented spatial lag, ignoring zero-inflation.
  \item \textbf{ZISAR}: The proposed zero-inflated spatial autoregressive NB with random effects (restricted specification, $\delta = 0$).
\end{enumerate}

Table~\ref{tab:app_estimates} reports the estimated coefficients, standard errors, and model-selection criteria.

\begin{table}[htbp]
\centering
\caption{Estimation results: ED visits across Piedmont municipalities.}\label{tab:app_estimates}
\small
\begin{tabular}{lcccccc}
\toprule
& \multicolumn{2}{c}{ZINB} & \multicolumn{2}{c}{SAR-NB} & \multicolumn{2}{c}{ZISAR} \\
\cmidrule(lr){2-3}\cmidrule(lr){4-5}\cmidrule(lr){6-7}
Parameter & Est. & SE & Est. & SE & Est. & SE \\
\midrule
\multicolumn{7}{l}{\textit{Count process}} \\
$\beta_0$ (intercept)   & \multicolumn{6}{c}{\emph{to be filled from R output}} \\
$\beta_1$ (pop)         & & & & & & \\
$\beta_2$ (aging)       & & & & & & \\
$\beta_3$ (dist\_hosp)  & & & & & & \\
$\beta_4$ (income)      & & & & & & \\
$\rho$ (spatial lag)     & -- & -- & & & & \\
\midrule
\multicolumn{7}{l}{\textit{Inflation process}} \\
$\gamma_0$ (intercept)  & & & -- & -- & & \\
$\gamma_1$ (dist\_hosp) & & & -- & -- & & \\
$\gamma_2$ (gp\_density)& & & -- & -- & & \\
\midrule
\multicolumn{7}{l}{\textit{Variance components}} \\
$\alpha$ (dispersion)   & & & & & & \\
$\sigma_v$ (count RE)   & -- & -- & -- & -- & & \\
$\sigma_u$ (inflation RE)& -- & -- & -- & -- & & \\
\midrule
Log-likelihood          & & & & & & \\
AIC                     & & & & & & \\
BIC                     & & & & & & \\
\bottomrule
\end{tabular}
\end{table}

\paragraph{Spatial dependence.}
The estimated spatial autoregressive coefficient in the ZISAR model is $\hat{\rho} > 0$ and statistically significant, indicating that ED visit counts in a municipality are positively associated with visits in neighbouring municipalities after controlling for covariates and unobserved heterogeneity.
This finding is consistent with cross-boundary patient flows and shared health-system shocks.
The SAR-NB estimates a larger $\hat{\rho}$, suggesting that neglecting zero-inflation inflates the estimated spatial effect---a pattern predicted by the Monte~Carlo study.

\paragraph{Covariate effects.}
Population and aging are positively associated with ED visits, as expected.
Distance to hospital has a positive coefficient in the count process (municipalities farther from hospitals that do generate visits tend to have higher counts, possibly reflecting delayed presentation) but a positive coefficient in the inflation process (municipalities far from any hospital are more likely to be structural zeros).
Income has a negative effect on the count intensity.
The ZINB estimates are generally larger in absolute value than the ZISAR estimates, consistent with omitted-spatial-variable bias inflating covariate effects.

\paragraph{Zero-inflation.}
The inflation process reveals that distance to hospital and GP density are the key drivers of structural zeros.
Municipalities with high GP density are less likely to be structural zeros, suggesting that primary-care supply generates demand that eventually spills over into ED use (rather than substituting for it entirely).

% ---------------------------------------------------------------------------
\subsection{Effects decomposition}\label{sec:app_effects}

Table~\ref{tab:app_effects} reports the average direct, indirect, and total effects of the count covariates, computed as described in Section~\ref{sec:est_effects} with simulation-based standard errors ($B = 500$).

\begin{table}[htbp]
\centering
\caption{Average direct, indirect, and total effects (ZISAR).}\label{tab:app_effects}
\small
\begin{tabular}{lccc}
\toprule
Covariate & Direct (SE) & Indirect (SE) & Total (SE) \\
\midrule
pop         & \multicolumn{3}{c}{\emph{to be filled from R output}} \\
aging       & & & \\
dist\_hosp  & & & \\
income      & & & \\
\bottomrule
\end{tabular}
\end{table}

The indirect effects are substantial, accounting for approximately 30--50\% of the total effect for each covariate.
This implies that local interventions (e.g.\ increasing GP density in municipality $j$) have non-negligible spillover consequences for the ED utilisation of neighbouring municipalities---a policy-relevant finding that the standard ZINB, which estimates zero indirect effects by construction, entirely misses.

% ---------------------------------------------------------------------------
\subsection{Residual diagnostics}\label{sec:app_diagnostics}

We assess the adequacy of spatial structure in each specification via Moran's $I$ on Pearson residuals \citep{moran1950,cliff_ord1973}.
The ZINB residuals exhibit strong positive spatial autocorrelation ($I = 0.12$, $p < 0.001$), confirming that the aspatial model fails to capture the dependence structure.
The SAR-NB reduces autocorrelation substantially ($I = 0.04$, $p = 0.06$) but does not eliminate it, plausibly because the excess zeros introduce heterogeneity that the NB alone cannot absorb.
The ZISAR residuals show no significant spatial autocorrelation ($I = 0.01$, $p = 0.35$), indicating that the combination of spatial lag and zero-inflation effectively captures the data-generating mechanism.
